\documentclass[11pt,a4paper]{article}

\usepackage[utf8]{inputenc}
\usepackage[T1]{fontenc}
\usepackage{geometry}
\usepackage{hyperref}
\usepackage{booktabs}
\usepackage{graphicx}
\usepackage{amsmath}
\usepackage{enumitem}
\usepackage{titlesec}
\usepackage{parskip}
\usepackage{cite}
\usepackage{url}
\usepackage{xcolor}
\usepackage{caption}
\usepackage{float}

\geometry{margin=2.5cm}

\hypersetup{
    colorlinks=true,
    linkcolor=blue,
    urlcolor=blue,
    citecolor=blue
}

\title{
    \textbf{CENG 467 --- Natural Language Understanding and Generation} \\[0.5em]
    \large Progress Report \\[0.5em]
    \LARGE Evaluating Context and History Pruning Strategies \\
    for RAG Faithfulness in Small Language Models \\
    via LLM-as-Judge
}

\author{
    Mert G\"{u}den (300201013) \and Berkay Fehmi Tekin (300201090)
}

\date{May 2026}

\begin{document}

\maketitle

\section{Project Title, Group Members, and Problem Definition}

This project investigates how context pruning strategies affect the faithfulness of answers generated by small language models (SLMs) in a retrieval-augmented generation (RAG) setting. Faithfulness---the degree to which a generated answer is grounded in the retrieved context---is a critical quality dimension for RAG systems, yet it remains difficult to measure automatically and is particularly vulnerable when the context contains noisy or irrelevant passages.

Our task is open-domain question answering over HotpotQA \cite{yang2018hotpotqa} (distractor setting), a multi-hop QA benchmark in which each question is accompanied by ten passages: two supporting passages and eight distractors drawn from Wikipedia. For each question, our pipeline retrieves the top-5 most relevant passages using BM25, optionally applies a pruning strategy to reduce context length and noise, generates a short answer using a small instruction-tuned language model (Mistral-7B-Instruct or Llama-3.1-8B-Instruct), and evaluates the answer's faithfulness using an LLM-as-Judge pipeline that decomposes the answer into atomic claims and verifies each claim against the retrieved context.

The central technical challenge is threefold. First, faithfulness is not directly measurable without expensive human annotation; we address this with a calibrated LLM-as-Judge \cite{es2023ragas} anchored to Cohen's $\kappa$ against human labels. Second, pruning introduces a coverage--noise trade-off: aggressive compression may remove the very sentences that support the correct answer. Third, small language models hallucinate more frequently under long or noisy contexts, making the pruning step crucial to generation quality. The seven pruning strategies we evaluate span from no pruning at all to combined semantic and history-aware compression, allowing us to systematically isolate the effect of each design choice.

\section{Dataset / Benchmark Status}

We use HotpotQA in its distractor configuration, loaded via the HuggingFace \texttt{datasets} library (\texttt{load\_dataset("hotpot\_qa", "distractor")}). The dataset contains approximately 90,447 training samples and 7,405 validation samples. For our experiments we draw 200 samples from the validation set and 50 separate samples for LLM-as-Judge calibration against human annotations. We do not fine-tune any model; all experimental samples come from the validation split.

Each sample provides a natural language question, a short gold answer (typically a span or a yes/no value), ten passages with their titles, and supporting-fact labels indicating which sentences are relevant. During preprocessing we flatten each passage's sentence list into a single string, apply whitespace tokenization for BM25 indexing, and truncate individual passages to a configurable token budget where needed. The loader script (\texttt{src/data/loader.py}) saves cleaned records in JSONL format. Dataset loading is fully implemented; the processed JSONL files are ready for the retrieval and pruning stages.

The primary challenges introduced by HotpotQA are: (1) the eight distractor passages directly inflate context length and introduce factual noise; (2) multi-hop questions require evidence from at least two separate passages, so aggressive pruning risks discarding one of the required supporting passages; and (3) HotpotQA does not provide a dedicated retrieval split, so we evaluate end-to-end retrieval alongside pruning on the same validation samples.

\begin{table}[H]
\centering
\caption{HotpotQA Dataset Summary}
\begin{tabular}{ll}
\toprule
\textbf{Property} & \textbf{Value} \\
\midrule
Full name & HotpotQA (distractor setting) \\
Source & HuggingFace \texttt{hotpot\_qa} \\
Training samples & $\sim$90,447 \\
Validation samples & $\sim$7,405 \\
Samples used (experiments) & 200 \\
Samples used (calibration) & 50 \\
Passages per question & 10 (2 supporting + 8 distractors) \\
Answer type & Short span / Yes--No \\
Task type & Multi-hop open-domain QA \\
Status & Loader implemented, JSONL files prepared \\
\bottomrule
\end{tabular}
\end{table}

\section{Literature Review Progress}

\paragraph{RAG and faithfulness.}
The RAG framework was introduced by Lewis et al.\ \cite{lewis2020rag} as a method for grounding language model outputs in dynamically retrieved documents, improving performance on knowledge-intensive tasks. Shuster et al.\ \cite{shuster2021retrieval} demonstrated empirically that retrieval augmentation substantially reduces hallucination in conversational settings compared to purely parametric generation. More recently, Min et al.\ \cite{min2023factscore} proposed FActScore, a fine-grained evaluation methodology that decomposes long-form generated text into atomic facts and verifies each against a knowledge source---an approach that directly motivates our LLM-as-Judge faithfulness metric.

\paragraph{Context compression.}
Jiang et al.\ \cite{jiang2023llmlingua} introduced LLMLingua, a token-level prompt compression approach that uses a small proxy language model to assign importance scores to individual tokens and removes low-scoring ones while preserving semantic integrity. Pan et al.\ \cite{pan2024llmlingua2} extended this line of work with LLMLingua-2, improving compression efficiency through data distillation and achieving better task-agnostic performance. Yoon et al.\ \cite{yoon2024provence} proposed Provence, a sentence-level pruning model trained specifically for RAG settings that scores each sentence's relevance to the query and discards below-threshold content. Xu et al.\ \cite{xu2023recomp} presented RECOMP, which frames context compression as either extractive or abstractive summarization, selecting query-relevant sentences via a learned compressor trained with contrastive objectives.

\paragraph{LLM-as-Judge.}
Zheng et al.\ \cite{zheng2023mtbench} systematically evaluated the feasibility of using strong LLMs as judges for open-ended generation quality, identifying positional and verbosity biases but demonstrating high correlation with human rankings on MT-Bench. Es et al.\ \cite{es2023ragas} introduced RAGAS, a reference-free framework for RAG evaluation that measures faithfulness by checking whether each claim in the generated answer is entailed by the retrieved context---the same decompose-then-verify paradigm we adopt. We calibrate our judge using Cohen's $\kappa$ against human annotations on 50 held-out samples, following standard NLP practice for inter-annotator agreement measurement.

\paragraph{Technical direction.}
Our project follows the RAGAS faithfulness pipeline but uses open-weight models rather than proprietary APIs. The baseline pruners (B1--B3) replicate the experimental setup from the RECOMP paper; the advanced pruners (S1--S4) are based on LLMLingua-2 and Provence with additional history-aware extensions that remove redundant content from prior conversational turns.

\section{Baseline Models and Pruning Strategies}

We evaluate seven pruning strategies organized into three baselines and four advanced methods. The baselines establish reference points that bracket expected performance: B1 (no pruning) represents the upper bound for context coverage but the worst noise level; B2 (naive truncation) is the standard production fallback that limits total context to 512 tokens by hard cutoff; B3 (RECOMP extractive) is the strongest prior extractive baseline, selecting query-relevant sentences via Jaccard token overlap. The four advanced strategies (S1--S4) are our proposed contributions. All strategies share the same \texttt{BasePruner} interface and are registered in a central \texttt{PRUNER\_REGISTRY}, enabling fair comparison under identical retrieval and generation conditions.

\begin{table}[H]
\centering
\caption{Baseline and Advanced Pruning Strategies}
\begin{tabular}{llll}
\toprule
\textbf{ID} & \textbf{Strategy} & \textbf{Type} & \textbf{Status} \\
\midrule
B1 & No Pruning (full context) & Baseline & \textbf{Complete} \\
B2 & Naive Truncation (512 tok) & Baseline & \textbf{Complete} \\
B3 & RECOMP Extractive & Prior work baseline & \textbf{Complete} \\
S1 & LLMLingua-2 & Token-level compression & \textbf{Complete} \\
S2 & Provence & Sentence-level pruning & \textbf{Complete} \\
S3 & Semantic History Pruning & History-aware pruning & \textbf{Complete} \\
S4 & Provence + History (combined) & Combined & \textbf{Complete} \\
\bottomrule
\end{tabular}
\end{table}

B1 serves as the upper-bound noise reference: providing the full retrieved context gives the generator the maximum available evidence but also exposes it to all distractor content. B2 is justified as the standard engineering fallback in production RAG systems where latency budgets constrain context length. B3 bridges the gap between simple truncation and learned compression by selecting sentences that lexically overlap with the query, serving as the extractive prior-work baseline. S1 (LLMLingua-2) and S2 (Provence) represent the current state of the art in token-level and sentence-level compression respectively. S3 (history pruning) addresses a limitation of static pruners by removing passages that overlap with information already surfaced in prior conversational turns. S4 (combined) chains Provence followed by history pruning, aiming to benefit from both semantic filtering and redundancy removal.


\section{Experimental Results}

All reported experiments were reproduced in Google Colab notebooks. Notebook~2 provides a baseline sanity run over 50 HotpotQA validation samples for the original seven-pruner comparison. Notebook~5 provides the final gap-closing runs on a held-out 50-sample slice (validation indices 600--649), specifically targeting the issues identified in the progress report: genuine LLMLingua-2 execution, non-empty two-turn history pruning, and qualitative error-analysis exports. In both notebooks, Mistral-7B-Instruct-v0.2 is used as the generator and as the LLM-as-Judge, with the loaded model shared between generation and judging when possible.

\subsection{Baseline Colab Run}

Table~\ref{tab:baseline-results} summarizes the Notebook~2 baseline run. This run is useful as the broad seven-pruner reference point and confirms the expected preliminary pattern: RECOMP improves faithfulness over hard truncation, while Provence's default threshold is very aggressive for multi-hop HotpotQA.

\begin{table}[H]
\centering
\caption{Notebook 2 baseline results --- 50 HotpotQA validation samples, Mistral-7B-Instruct-v0.2}
\label{tab:baseline-results}
\begin{tabular}{lccccr}
\toprule
\textbf{Strategy} & \textbf{Faithfulness} & \textbf{EM} & \textbf{F1} & \textbf{Comp. Ratio} & \textbf{Lat. (s)} \\
\midrule
B1: No Pruning       & 0.718 & 0.040 & 0.148 & 1.000 & 6.76 \\
B2: Naive Truncation & 0.621 & 0.040 & 0.139 & 0.918 & 4.95 \\
B3: RECOMP           & \textbf{0.750} & 0.040 & 0.141 & 0.911 & 5.38 \\
S1: LLMLingua-2 (fallback) & \textbf{0.750} & 0.040 & 0.141 & 0.911 & 5.38 \\
S2: Provence         & 0.657 & 0.000 & 0.057 & \textbf{0.072} & \textbf{4.49} \\
S3: History (empty) & 0.718 & 0.040 & 0.148 & 1.000 & 6.34 \\
S4: Combined (empty) & 0.657 & 0.000 & 0.057 & \textbf{0.072} & \textbf{4.48} \\
\bottomrule
\end{tabular}
\end{table}

\noindent\textit{Note:} The S1 row in Table~\ref{tab:baseline-results} is retained only as a historical baseline artifact: in this run LLMLingua-2 fell back to the RECOMP implementation, which explains the identical S1/B3 scores. Similarly, S3 and S4 use empty history and therefore represent degenerate single-turn cases.

\subsection{Final Gap-Closing Colab Run}

Notebook~5 re-ran the problematic strategies under corrected conditions. LLMLingua-2 was executed with the \texttt{llmlingua} package and \texttt{use\_llmlingua2=True}; the recorded fallback counter is zero. History pruning and combined pruning were evaluated with synthetic two-turn history rather than empty history. Table~\ref{tab:gap-results} reports these results, including supporting-fact coverage and noise ratio to make pruning failures visible beyond faithfulness alone.

\begin{table}[H]
\centering
\caption{Notebook 5 final gap-closing results --- 50 held-out HotpotQA validation samples (indices 600--649)}
\label{tab:gap-results}
\begin{tabular}{lcccccc}
\toprule
\textbf{Strategy} & \textbf{Faith.} & \textbf{EM} & \textbf{F1} & \textbf{Comp.} & \textbf{Coverage} & \textbf{Noise} \\
\midrule
B3: RECOMP & 0.648 & 0.080 & 0.195 & 0.962 & 0.760 & 0.624 \\
S1: LLMLingua-2 (real) & \textbf{0.711} & 0.040 & 0.167 & 0.457 & 0.000 & 0.792 \\
S3: History (two-turn) & 0.533 & \textbf{0.100} & \textbf{0.214} & 0.962 & 0.627 & 0.641 \\
S4: Combined (two-turn) & 0.530 & 0.040 & 0.112 & \textbf{0.081} & 0.222 & \textbf{0.284} \\
\bottomrule
\end{tabular}
\end{table}

The corrected LLMLingua-2 result is no longer identical to RECOMP: it compresses much more aggressively (0.457 vs. 0.962) and obtains higher judge faithfulness (0.711 vs. 0.648), but its supporting-fact coverage metric is 0.000 under our exact/fuzzy retained-sentence check. This indicates that LLMLingua-2 often rewrites or fragments evidence in a way that can still support model answers but no longer preserves gold supporting sentences verbatim. Two-turn history pruning does change behaviour relative to the empty-history case, but on this slice it reduces faithfulness to 0.533, suggesting that the current novelty-only redundancy score can discard useful repeated evidence. The combined Provence+History setting is the most compact (0.081) and lowest-noise (0.284), but the low coverage (0.222) confirms that the current Provence threshold remains too aggressive for multi-hop QA.

\section{Completed Improvements and Technical Direction}

\paragraph{Completed.}
The final Colab workflow now addresses three of the major report weaknesses. First, the GitHub repository link is fixed and the report references the notebooks rather than local Python commands. Second, LLMLingua-2 is evaluated as a genuine LLMLingua-2 run with zero fallback cases. Third, S3 and S4 are evaluated with synthetic two-turn history, so history pruning is no longer reported only in the single-turn degenerate setting. Notebook~5 also exports \texttt{error\_analysis\_candidates.jsonl}, which provides concrete examples for qualitative analysis.

\paragraph{Remaining.}
LLM-as-Judge calibration is scaffolded but not yet fully completed in the checked-in artifacts: Notebook~5 creates the annotation template and calibration computation cells, but no filled \texttt{human\_labels.jsonl} or resulting Cohen's $\kappa$ file is present in the repository output folders. This should be treated as remaining work rather than claimed as completed. Provence threshold tuning also remains important because both S2 and S4 show very low compression ratios and supporting-fact coverage under the default threshold.

\paragraph{Future extension.}
The LNN/CfC direction is not yet a completed experimental contribution. We implemented exploratory Colab notebooks for LNN-style sentence/token pruning, but the training and evaluation loop has not reached a stable enough point to support quantitative claims in this report. We therefore position Liquid Neural Networks and Closed-form Continuous-time (CfC) models as future adaptive pruning controllers. The intended role is to choose pruning aggressiveness dynamically from query complexity, retrieval redundancy, and history overlap, instead of relying on a fixed Provence threshold or a purely lexical history novelty score. This direction is especially relevant for multi-turn QA, where evidence importance changes across turns.

A second future extension is dataset breadth. HotpotQA is useful for multi-hop distractor QA, but it does not fully represent conversational retrieval behaviour. Future experiments should add conversational QA benchmarks such as QuAC and QReCC-style query rewriting / conversational context datasets, where history pruning can be evaluated against naturally occurring dialogue context rather than synthetic two-turn histories.

\section{Ablation and Prompt Sensitivity Plan}
\label{sec:ablation}

We plan a targeted set of ablation experiments to isolate the contribution of individual design choices:

\begin{itemize}[nosep]
    \item \textbf{Judge prompt format:} zero-shot vs.\ 2-shot examples for both claim decomposition and claim verification steps.
    \item \textbf{Verification response type:} strict binary YES/NO vs.\ soft confidence score thresholded at 0.5.
    \item \textbf{Retrieval depth:} $k=3$ vs.\ $k=5$ retrieved passages (affects both coverage and noise level).
    \item \textbf{Truncation budget:} \texttt{max\_tokens} = 256 vs.\ 512 for B2 and B3.
    \item \textbf{History context:} S3 and S4 evaluated with empty history (single-turn) vs.\ a simulated two-turn history to assess the incremental value of history-aware pruning.
\end{itemize}

These ablations are designed to answer whether the faithfulness gains from advanced pruners are attributable to the compression strategy itself or to incidental factors such as prompt wording or retrieval depth.


\section{Error Analysis and Generation Quality}

Notebook~5 exports concrete error-analysis candidates rather than only a category plan. The observed errors fall into four recurring groups.

\paragraph{Pruning removes supporting evidence.}
For the question ``The song Arizona was recorded by Paul Revere and Mark Lindsay but who wrote the song?'', the gold answer is \textit{Kenny Young}. Under the real LLMLingua-2 run, the generated answer is ``I don't know,'' and supporting-fact coverage is 0.000. The exported supporting sentence explicitly states that ``Arizona'' was written by Kenny Young, so this is a clear pruning/coverage failure. RECOMP and two-turn history pruning also partially lose evidence on the same example (coverage 0.5), although RECOMP still produces the correct entity.

\paragraph{Exact match under-credits verbose but correct answers.}
For the yes/no question ``Are both Parodia and Thalictrum flowering plants?'', S4 generates: ``Yes, both Parodia and Thalictrum are flowering plants,'' followed by a short explanation. Faithfulness is 1.0 and coverage is 1.0, but EM is 0.0 because the gold answer is simply \textit{yes}. This confirms that EM is a poor primary metric for instruction-tuned generators that produce explanatory answers.

\paragraph{Judge failures occur on short supported answers.}
For the Watercliffe Meadow question, the history-pruning run generates exactly \textit{Political correctness}; EM and F1 are both 1.0, coverage is 1.0, but the judge faithfulness score is 0.0. This is likely an LLM-as-Judge false negative caused by the answer being a short phrase rather than a full claim. This case motivates the remaining calibration work.

\paragraph{Near-miss entity generation.}
For the ``Black Maverick'' biography question, RECOMP retrieves evidence for \textit{T. R. M. Howard}, but the generator outputs \textit{T. M. Howard}. Token F1 remains high (0.857), but EM is 0.0 and faithfulness is 0.0. This illustrates a generation-level entity omission error: the context contains the answer, but the generated string drops an initial.

Overall, the strongest pattern is that faithfulness, EM/F1, and supporting-fact coverage measure different failure modes. LLMLingua-2 obtains the best faithfulness in the final gap-closing run, but its zero retained-sentence coverage indicates that the coverage metric is too strict for compressed/rephrased contexts or that LLMLingua-2 is removing gold evidence while retaining enough related context for the generator. Conversely, some exact or semantically correct answers receive low judge scores, so calibration remains necessary before treating judge faithfulness as a standalone ground truth.

\section{Ethical Considerations and Bias}

The LLM-as-Judge pipeline is powered by Mistral-7B-Instruct, a model that may inherit the biases present in its pretraining and instruction-tuning corpora, including positional bias (preferring claims that appear early in the context), verbosity bias (favoring longer answers as more faithful), and cultural or demographic skews. HotpotQA passages are sourced from English Wikipedia, which is known to underrepresent non-Western entities and perspectives; our faithfulness measurements may therefore be systematically more reliable for topics with dense Wikipedia coverage. We also note that faithfulness to retrieved context is not equivalent to truthfulness about the world: a model that correctly reproduces a false claim from a retrieved passage will score high on faithfulness while being factually incorrect. All models used in this project are open-weight and used solely for academic research; no proprietary API access is required, and no user data or personally identifiable information is involved. Results from this study should not be extrapolated to high-stakes QA applications without additional human evaluation.

\section{GitHub and Reproducibility Status}

The repository is organized as a modular Python package with clearly separated components for data loading, retrieval, pruning, generation, evaluation, and judge calibration. For reproduction, however, the primary execution interface is Google Colab rather than direct local Python commands. The final report will reference the Colab notebooks as the runnable artifacts: \texttt{01\_data\_exploration.ipynb} for dataset checks, \texttt{02\_baseline\_results.ipynb} for small-scale baseline sanity checks, and \texttt{05\_final\_colab\_runs.ipynb} for the final LLMLingua-2, two-turn history, judge calibration, and error-analysis runs. The LNN/CfC notebooks are kept as exploratory artifacts rather than completed result notebooks.

\begin{verbatim}
rag-faithfulness-pruning/
|-- PROJECT.md
|-- README.md
|-- requirements.txt
|-- test_pipeline.py
|-- notebooks/
|   |-- 01_data_exploration.ipynb
|   |-- 02_baseline_results.ipynb
|   |-- 03_lnn_training_colab.ipynb
|   |-- 04_lnn_evaluation_colab.ipynb
|   \-- 05_final_colab_runs.ipynb
|-- data/
|   |-- raw/          (hotpotqa_train.jsonl, hotpotqa_validation.jsonl)
|   |-- processed/    (hotpotqa_experiments.jsonl, hotpotqa_calibration.jsonl)
|   \-- judge_calibration/
|-- src/
|   |-- data/         (loader.py, splitter.py)
|   |-- retrieval/    (bm25_retriever.py)
|   |-- pruning/      (base.py, no_pruning.py, naive_truncation.py,
|   |                  recomp.py, llmlingua2.py, provence.py,
|   |                  history_pruning.py, combined.py)
|   |-- generation/   (generator.py)
|   |-- judge/        (judge.py, prompts.py, calibration.py)
|   \-- evaluation/   (runner.py, metrics.py)
\-- experiments/results/
    |-- notebook2/   (baseline aggregate summary and sanity outputs)
    \-- notebook5/   (final gap-closing JSONL files, summary, error cases)
\end{verbatim}

All seven core pruners are implemented and available through the shared registry. Final reported outputs are reproduced by opening \texttt{notebooks/05\_final\_colab\_runs.ipynb} in a GPU Colab runtime and running the cells from top to bottom. The notebook clones the repository, installs dependencies, downloads HotpotQA, writes JSONL result files under \texttt{experiments/results/}, produces a report-ready aggregate summary, and exports candidate cases for qualitative error analysis.

GitHub repository: \url{https://github.com/Mrtuzy/CENG467_Final}


\section{Current Challenges and Next Steps}

\paragraph{Challenges.}
The final Colab results surface three concrete challenges. First, aggressive compression still creates a coverage--faithfulness tension. Real LLMLingua-2 obtains the highest faithfulness in the gap-closing run (0.711) with a much lower compression ratio (0.457), but the supporting-fact coverage metric drops to 0.000, suggesting either evidence loss or a mismatch between sentence-preservation metrics and token-level compressed contexts. Second, Provence-based compression remains too aggressive under the default threshold: the combined two-turn run keeps only 8.1\% of tokens and has low coverage (0.222), even though its retained context has the lowest noise ratio. Third, the judge still produces questionable scores on some short answers, including cases where EM/F1 are perfect but faithfulness is 0.0. This makes human calibration the most important remaining validation step.

\paragraph{Next steps.}
\begin{enumerate}[nosep]
    \item Complete LLM-as-Judge calibration by filling \texttt{human\_labels.jsonl} for at least 20--50 examples and reporting Cohen's $\kappa$.
    \item Tune Provence's threshold (e.g., 0.3 and 0.4) and re-run S2/S4 to target a less destructive compression ratio around 0.4--0.6.
    \item Add a compression-aware LLMLingua-2 coverage metric; exact retained-sentence coverage is too brittle for token-level compression.
    \item Scale the corrected Notebook~5 configuration from 50 to 200 samples after calibration and threshold tuning.
    \item Improve answer formatting by prompting the generator for a concise final answer, then report EM/F1 both before and after answer extraction.
    \item Stabilize the exploratory LNN/CfC notebooks before making quantitative claims; the next target is an adaptive pruning controller that chooses compression strength from query complexity, retrieval redundancy, and history overlap.
    \item Extend evaluation beyond HotpotQA to conversational QA datasets such as QuAC and QReCC-style benchmarks, where history pruning can be tested on natural dialogue histories rather than synthetic two-turn context.
\end{enumerate}

\begin{thebibliography}{10}

\bibitem{lewis2020rag}
P.\ Lewis, E.\ Perez, A.\ Piktus, F.\ Petroni, V.\ Karpukhin, N.\ Goyal, H.\ K\"{u}ttler, M.\ Lewis, W.\ Yih, T.\ Rockt\"{a}schel, S.\ Riedel, and D.\ Kiela.
\newblock Retrieval-Augmented Generation for Knowledge-Intensive NLP Tasks.
\newblock In \textit{Advances in Neural Information Processing Systems (NeurIPS)}, 2020.

\bibitem{shuster2021retrieval}
K.\ Shuster, S.\ Poff, M.\ Chen, D.\ Kiela, and J.\ Weston.
\newblock Retrieval Augmentation Reduces Hallucination in Conversation.
\newblock In \textit{Findings of EMNLP}, 2021.

\bibitem{xu2023recomp}
F.\ Xu, W.\ Shi, and E.\ Choi.
\newblock RECOMP: Improving Retrieval-Augmented LMs with Compression and Selective Augmentation.
\newblock In \textit{International Conference on Learning Representations (ICLR)}, 2024.

\bibitem{jiang2023llmlingua}
H.\ Jiang, Q.\ Wu, C.-Y.\ Lin, Y.\ Yang, and L.\ Qiu.
\newblock LLMLingua: Compressing Prompts for Accelerated Inference of Large Language Models.
\newblock In \textit{Proceedings of EMNLP}, 2023.

\bibitem{pan2024llmlingua2}
Z.\ Pan, Q.\ Wu, H.\ Jiang, M.\ Garg, N.\ Liden, Y.\ Ronen, and L.\ Qiu.
\newblock LLMLingua-2: Data Distillation for Efficient and Faithful Task-Agnostic Prompt Compression.
\newblock In \textit{Findings of ACL}, 2024.

\bibitem{yoon2024provence}
S.\ Yoon, W.\ Xu, F.\ Xu, and H.\ Ji.
\newblock Provence: Efficient and Robust Context Pruning for Retrieval-Augmented Generation.
\newblock In \textit{Proceedings of NAACL}, 2024.

\bibitem{zheng2023mtbench}
L.\ Zheng, W.-L.\ Chiang, Y.\ Sheng, S.\ Zhuang, Z.\ Wu, Y.\ Zhuang, Z.\ Lin, Z.\ Li, D.\ Li, E.\ P.\ Xing, H.\ Zhang, J.\ E.\ Gonzalez, and I.\ Stoica.
\newblock Judging LLM-as-a-Judge with MT-Bench and Chatbot Arena.
\newblock In \textit{Advances in Neural Information Processing Systems (NeurIPS)}, 2023.

\bibitem{es2023ragas}
S.\ Es, J.\ James, L.\ E.\ Anke, and S.\ Schockaert.
\newblock RAGAS: Automated Evaluation of Retrieval Augmented Generation.
\newblock In \textit{Proceedings of EACL}, 2024.

\bibitem{min2023factscore}
S.\ Min, K.\ Krishna, X.\ Lyu, M.\ Lewis, W.\ Yih, P.\ W.\ Koh, M.\ Iyyer, L.\ Zettlemoyer, and H.\ Hajishirzi.
\newblock FActScore: Fine-grained Atomic Evaluation of Factual Precision in Long Form Text Generation.
\newblock In \textit{Proceedings of EMNLP}, 2023.

\bibitem{yang2018hotpotqa}
Z.\ Yang, P.\ Qi, S.\ Zhang, Y.\ Bengio, W.\ W.\ Cohen, R.\ Salakhutdinov, and C.\ D.\ Manning.
\newblock HotpotQA: A Dataset for Diverse, Explainable Multi-hop Question Answering.
\newblock In \textit{Proceedings of EMNLP}, 2018.

\end{thebibliography}

\end{document}
